\noindent
\textbf{CS4268 AY2025-26 Sem 2}
\\
\noindent
Lecturer: Divesh Aggarwal
\hfill
Lecture 5                      %%% FILL IN LECTURE NUMBER HERE
\hfill
Date: February 12, 2026          %%% FILL IN LECTURE DATE HERE


\noindent
\rule{\textwidth}{1pt}

\medskip

\section{Basic Primitives in Quantum Computing}

In this section we introduce several fundamental techniques that will be used repeatedly throughout this text. They are uncomputation (which allows us to remove garbage and reuse ancilla qubits), the phase-kickback method (which encodes function values into phases), and the Hadamard-Walsh transform, which is a fancy name for the parallel implementation of
Hadamard gates.

\subsection{Uncomputation}

As discussed previously, a quantum implementation of a classical function
typically produces additional intermediate data, usually referred to as
garbage. Since quantum circuits are reversible, this garbage cannot be
simply discarded. However, we would like to reuse ancillas which requires them to be reverted back to their original state.

Let $F : \{0,1\}^n \rightarrow \{0,1\}$ and suppose we have a quantum circuit
$Q_F$ which, on input $\ket{x}$ and some ancilla qubits, produces
$\ket{F(x)}$ together with some garbage.

The standard technique to remove garbage is the
\textit{compute--copy--uncompute} sequence. We first compute $F(x)$,
then copy the result into a clean qubit using a CNOT, and finally apply the
inverse circuit $Q_F^{-1}$ to undo the computation and restore the ancillas. This process is commonly called \textbf{uncomputation}.
\[
\begin{quantikz}[row sep=0.7cm, column sep=0.7cm]
\lstick{$\ket{x}$} 
    & \gate[wires=2]{Q_F}
    & \ctrl{2} 
    & \gate[wires=2]{Q_F^{-1}} 
    & \rstick{$\ket{x}$} \qw \\
%
\lstick{$\ket{\text{ancillas}}$} 
    & \qw 
    & \qw 
    & \qw 
    & \rstick{$\ket{\text{ancillas}}$} \qw \\
%
\lstick{$\ket{0}$} 
    & \qw 
    & \targ{} 
    & \qw 
    & \rstick{$\ket{F(x)}$} \qw
\end{quantikz}
\]

After applying $Q_F^{-1}$ the garbage is removed and the ancillas are restored to their original state. We denote this entire implementation by $Q_F^* := Q_F^{-1} \circ \text{CNOT} \circ Q_F $.

More generally, for $F : \{0,1\}^n \rightarrow \{0,1\}^m$ and $y \in \{0,1\}^m$, the circuit $Q_F^*$ implements
\[
\begin{quantikz}
    \lstick{$\ket{x}$} & \gate[wires=2]{Q_F^*} & \rstick{$\ket{x}$} \\
    \lstick{$\ket{y}$} & \qw & \rstick{$\ket{y \oplus F(x)}$}
\end{quantikz}
\]
This form is essential because it avoids unwanted entanglement between the
output and the ancillas and allows the ancilla qubits to be reused.

\subsection{The Phase-Kickback Method}

In many quantum algorithms we work with Boolean functions $F : \{0,1\}^n \rightarrow \{0,1\}$. It is often convenient to use the equivalent phase representation
\[
f(x) = (-1)^{F(x)} \in \{-1,1\}.
\]

A key technique for obtaining this phase representation is the \textbf{phase-kickback} trick. Consider the  aforementioned $Q_F^*$. If we apply $Q_F^*$ to $\ket{x}$ and $\ket{-}$ we obtain $\frac{1}{\sqrt{2}} (\ket{F(x)} - \ket{\neg F(x)})$ as the second output which by closer inspection is the same as
\[
(-1)^{F(x)} \ket{x}\ket{-}.
\]
Thus the value of $F(x)$ is not stored in the target qubit; instead it appears
as a phase factor on $\ket{x}$. This is the phase kickback effect.

Using standard gates, we can prepare $\ket{-}$ from $\ket{0}$ by applying a
NOT gate followed by a Hadamard. This yields the circuit

\[
\begin{quantikz}
    \lstick{$\ket{x}$} & \qw & \qw & \gate[wires=2]{Q_F^*} &\qw & \qw & \rstick{$(-1)^{F(x)}\ket{x}$} \\
    \lstick{$\ket{0}$} & \gate{NOT} & \gate{H} & \qw & 
    \gate{H} & \gate{NOT} &
    \rstick{$\ket{0}$}
\end{quantikz}
\]

We denote this entire quantum circuit by $Q_F^{\pm}$.

\paragraph{Why do we do this?} Quantum algorithms exploit interference. Relative phases between basis states affect interference patterns. By encoding $F(x)$ as a phase rather than as a separate output bit, we allow the function values to influence interference directly.

\subsection{The Hadamard-Walsh Transform}

The operator $H^{\otimes n}$ denotes the parallel application of Hadamard gates to $n$ qubits. This transformation, also known as the
\textbf{Hadamard--Walsh transform}, maps
\[
\ket{0^n} \mapsto
\frac{1}{\sqrt{N}} \sum_{x \in \{0,1\}^n} \ket{x},
\qquad N = 2^n,
\]
creating a uniform superposition over all bit strings.
More generally,
\[
H^{\otimes n}\ket{x}
=
\frac{1}{\sqrt{N}}
\sum_{y \in \{0,1\}^n}
(-1)^{x \cdot y}\ket{y},
\]
where $x \cdot y$ is the bitwise inner product modulo $2$.

\paragraph{Example.}
To build intuition, consider the case $\ket{b_1,b_2}$.
Applying $H^{\otimes 2}$ gives
\[
\frac{1}{2}\Bigl(
\ket{00}
+ (-1)^{b_2}\ket{01}
+ (-1)^{b_1}\ket{10}
+ (-1)^{b_1+b_2}\ket{11}
\Bigr).
\]
Thus the output is a uniform superposition over all basis states, and the
input bits $b_1,b_2$ only determine the \emph{sign pattern} of the
amplitudes. The phases are given by the parities
$b_1x_1 + b_2x_2 \pmod{2}$.

In general, the application of $H^{\otimes n}$ maps a computational basis
state $\ket{x}$ to a vector with entries $\pm 1/\sqrt{2^n}$, where the
$s$-th coordinate is given by
\[
\frac{1}{\sqrt{2^n}}(-1)^{\sum_{i=1}^n x_i s_i \bmod 2}.
\]
where $s \in \{0,1\}^n$ labels the computational basis states.


These vectors form the so-called $\chi$-basis (the Fourier basis over
$\mathbb{Z}_2^n$). From this perspective, the Hadamard--Walsh transform is
simply the change of basis from the computational basis to the $\chi$-basis.
